% ABOUTME: The instruction subset and the single-cycle datapath and control unit.
% ABOUTME: Owner: Aryan (review against rtl/common and rtl/single_cycle, then defend in viva).
\section{Instruction Set and Single-Cycle Design}
\label{sec:single}

\subsection{Instruction subset}
The brief suggests ADD, SUB, ADDI, LW, SW, BEQ and BNE. We add the logic,
compare and shift instructions needed for realistic benchmarks such as
sorting, 18 instructions in all (Table~\ref{tab:isa}). Memory accesses are
aligned 32 bit words. \texttt{ecall} marks the end of a program.

\begin{table}[t]
\centering
\caption{Supported RV32I instructions}
\label{tab:isa}
\begin{tabular}{@{}ll@{}}
\toprule
Format & Instructions \\
\midrule
R & \texttt{add sub and or xor slt sll} \\
I & \texttt{addi andi ori xori slti slli} \\
I (load), S & \texttt{lw}, \texttt{sw} \\
B & \texttt{beq bne} \\
System & \texttt{ecall} (halt) \\
\bottomrule
\end{tabular}
\end{table}

\subsection{Datapath}
The datapath is built from shared modules in \texttt{rtl/common}: a 32 bit
ALU with seven operations, a $32\times32$ register file with two
asynchronous read ports and one synchronous write port (x0 always reads
zero), an immediate generator for the I, S and B formats, and two 256 word
memories. In the single-cycle CPU (\texttt{cpu\_single.v}) one instruction
passes through all of them in one clock period: the PC addresses instruction
memory, the fields address the register file, the ALU computes a result or an
address, data memory is read or written, and the write-back mux returns the
ALU result or the loaded word to the register file. A taken branch loads
$\mathrm{PC}+\mathrm{imm}$ instead of $\mathrm{PC}+4$.

\subsection{Control unit}
The control unit (\texttt{control.v}) decodes the opcode, \texttt{funct3}
and \texttt{funct7} into the signals of Table~\ref{tab:control}. It also
reports whether the instruction really reads \texttt{rs1} and \texttt{rs2}.
In an I-type instruction the \texttt{rs2} field holds immediate bits, and
without this flag the pipeline's hazard unit would stall on registers the
instruction never reads.

\begin{table}[t]
\centering
\caption{Main control signals per instruction class}
\label{tab:control}
\small
\setlength{\tabcolsep}{3pt}
\begin{tabular}{@{}lcccccc@{}}
\toprule
Class & RegWr & ALUSrc & MemRd & MemWr & Br & ALU \\
\midrule
R-type & 1 & reg & 0 & 0 & 0 & funct \\
I-type & 1 & imm & 0 & 0 & 0 & funct3 \\
\texttt{lw} & 1 & imm & 1 & 0 & 0 & add \\
\texttt{sw} & 0 & imm & 0 & 1 & 0 & add \\
\texttt{beq}/\texttt{bne} & 0 & reg & 0 & 0 & 1 & compare \\
\bottomrule
\end{tabular}
\end{table}

\subsection{Critical path}
Place and route on the ECP5 (Section~\ref{sec:results}) reports a longest
path of 17.0\,ns, or 58.75\,MHz. Because the instruction memory is a ROM built
from logic, Yosys moves the PC register to the ROM's output, so the path starts
with the instruction already in flip-flops. From there it runs through a
register file read (2.1\,ns), the ALU computing the address (3.2\,ns), a read of
the 256 word data memory (8.0\,ns, the largest share) and the write-back mux to
the register file's write port (3.2\,ns, plus setup). Every instruction pays
for this path, even an \texttt{addi} that never touches data memory.
